% Options for packages loaded elsewhere
\PassOptionsToPackage{unicode}{hyperref}
\PassOptionsToPackage{hyphens}{url}
%
\documentclass[
]{article}
\usepackage{amsmath,amssymb}
\usepackage{iftex}
\ifPDFTeX
  \usepackage[T1]{fontenc}
  \usepackage[utf8]{inputenc}
  \usepackage{textcomp} % provide euro and other symbols
\else % if luatex or xetex
  \usepackage{unicode-math} % this also loads fontspec
  \defaultfontfeatures{Scale=MatchLowercase}
  \defaultfontfeatures[\rmfamily]{Ligatures=TeX,Scale=1}
\fi
\usepackage{lmodern}
\ifPDFTeX\else
  % xetex/luatex font selection
\fi
% Use upquote if available, for straight quotes in verbatim environments
\IfFileExists{upquote.sty}{\usepackage{upquote}}{}
\IfFileExists{microtype.sty}{% use microtype if available
  \usepackage[]{microtype}
  \UseMicrotypeSet[protrusion]{basicmath} % disable protrusion for tt fonts
}{}
\makeatletter
\@ifundefined{KOMAClassName}{% if non-KOMA class
  \IfFileExists{parskip.sty}{%
    \usepackage{parskip}
  }{% else
    \setlength{\parindent}{0pt}
    \setlength{\parskip}{6pt plus 2pt minus 1pt}}
}{% if KOMA class
  \KOMAoptions{parskip=half}}
\makeatother
\usepackage{xcolor}
\setlength{\emergencystretch}{3em} % prevent overfull lines
\providecommand{\tightlist}{%
  \setlength{\itemsep}{0pt}\setlength{\parskip}{0pt}}
\setcounter{secnumdepth}{-\maxdimen} % remove section numbering
\ifLuaTeX
  \usepackage{selnolig}  % disable illegal ligatures
\fi
\usepackage{bookmark}
\IfFileExists{xurl.sty}{\usepackage{xurl}}{} % add URL line breaks if available
\urlstyle{same}
\hypersetup{
  hidelinks,
  pdfcreator={LaTeX via pandoc}}

\author{}
\date{}

\begin{document}

\section{A bounded weak ray with trivial non-Archimedean datum and
infinite entropy
slope}\label{a-bounded-weak-ray-with-trivial-non-archimedean-datum-and-infinite-entropy-slope}

\textbf{Status.} Root-authored candidate, not independently reviewed.
This proves an obstruction to the unrestricted entropy-slope identity.
The original P19.1 residual problem remains open pending an independent
comparison with its exact completion criterion. In particular, an
example does not establish the entropy-slope theorem for every maximal
ray on every compact Kähler manifold.

\subsection{1. Precise statement and
normalizations}\label{precise-statement-and-normalizations}

Let \((Y,\eta)\) be any compact Kähler manifold of dimension \(m\geq0\),
with the point allowed when \(m=0\). Put \[
 X=\mathbf P^1\times Y,\qquad
 \Theta=p_1^*\omega+p_2^*\eta,\qquad
 V_X=\int_X\Theta^{m+1},
\] where \(\omega\) is the Fubini--Study form of mass one. We use
\(dd^c=(i/\pi)\partial\bar\partial\), so that \(dd^c\log|z|=\delta_0\)
on a complex curve. For bounded \(\Theta\)-psh \(u\), use the normalized
energy and probability measure \[
 E_\Theta(u)=\frac{1}{(m+2)V_X}
 \sum_{j=0}^{m+1}\int_Xu\,\Theta_u^j\wedge\Theta^{m+1-j},
 \qquad \operatorname{MA}_\Theta(u)=V_X^{-1}\Theta_u^{m+1}.
\] Set \(\mu_0=V_X^{-1}\Theta^{m+1}\). The entropy is relative entropy
of probability measures: \[
 \operatorname{Ent}_\Theta(u)=
 \begin{cases}\int_X\rho\log\rho\,d\mu_0,&\operatorname{MA}_\Theta(u)=\rho\mu_0,\\
 +\infty,&\operatorname{MA}_\Theta(u)\not\ll\mu_0.
 \end{cases}
\] The ray time is \(t=-\log|\zeta|\). If a vertical prime divisor \(F\)
on a modification of \(X\times\mathbf P^1\) has
\(b_F=\operatorname{ord}_F(\zeta)\), its normalized valuation on ideal
sheaves from \(X\) is \(v_F=b_F^{-1}\operatorname{ord}_F\). We use the
datum \[
 u^{\rm NA}(v_F)=-b_F^{-1}\nu_F(U),\qquad U(x,\zeta)=u_{-\log|\zeta|}(x),
\] where \(\nu_F(U)\) is the generic divisorial coefficient of its
pullback, in the convention \(dd^c\log|z_1|=[z_1=0]\). Divisorial
determination fixes its psh extension. The discrepancy for entropy is
\textbf{over \(X\)}:
\(A_X(c\operatorname{ord}_D)=c(1+\operatorname{ord}_D K_{Z/X})\) on a
smooth modification \(Z\to X\), extended in the usual way, and
\(A_X(v_{\rm triv})=0\). Thus the discrepancy of the Gauss extension
over \(X\times\mathbf P^1\) is shifted by its time weight before it
enters entropy.

\textbf{Candidate theorem.} For every \(a\in(0,1)\) there is a bounded
weak geodesic ray \((u_t)_{t\geq0}\) from zero such that \[
 -at\leq u_t\leq0,\qquad E_\Theta(u_t)=-\frac{a^2}{2}t,
 \qquad u^{\rm NA}=0,
 \qquad \lim_{t\to\infty}\frac{\operatorname{Ent}_\Theta(u_t)}{t}=+\infty.
 \tag{1}
\] Here bounded means bounded at every finite time; no time-independent
bound on \(\|u_t\|_\infty\) is asserted. The non-Archimedean metric zero
has normalized Monge--Ampère measure \(\delta_{v_{\rm triv}}\), and
hence \[
 H^{\rm NA}(u^{\rm NA})=\int_{X^{\beth}} A_X\,d\operatorname{MA}^{\rm NA}_\Theta(0)=0.
 \tag{2}
\] Thus the ordinary non-Archimedean datum does not determine entropy
slopes of all bounded weak rays: the zero ray has the same datum and
entropy slope zero, whereas (1) has infinite entropy slope. Taking \(Y\)
nonprojective gives a compact nonprojective example. Taking an arbitrary
transcendental Kähler class on \(Y\) gives the corresponding example in
a transcendental product class.

\subsection{2. An explicit atom-free polar probability
measure}\label{an-explicit-atom-free-polar-probability-measure}

Start with \(I_0=[0,1]\subset\mathbf C\subset\mathbf P^1\). For
\(j\geq1\), let \(\delta_j=\exp(-4^j)\). Replace every interval at level
\(j-1\) by its two closed end intervals of length \(\delta_j\). These
intervals are disjoint: for \(j\geq2\), \(2\delta_j<\delta_{j-1}\), and
\(2\delta_1<1\). Let \(K\) be the intersection of these decreasing
unions of intervals.

The fair Bernoulli measure on the binary choices pushes forward to a
probability measure \(\sigma\) on \(K\). Every level-\(j\) interval has
mass \(2^{-j}\). A point belongs to a unique interval at each level, so
\(\sigma(\{x\})\leq2^{-j}\) for every \(j\), and \(\sigma\) has no
atoms.

On the affine chart write the Green kernel \[
 G(z,w)=\log\frac{|z-w|}{\sqrt{1+|z|^2}\sqrt{1+|w|^2}}\leq0.
\] Its extension across infinity satisfies
\(dd_z^cG(z,w)=\delta_w-\omega\). Define \[
 \psi(z)=\int_KG(z,w)\,d\sigma(w).
 \tag{3}
\] The local integrability of the logarithmic kernel, uniformly for
\(w\) in the compact interval \([0,1]\), and Fubini's theorem show that
\(\psi\) is not identically \(-\infty\) and belongs to \(L^1(\omega)\).
Distributionally, \[
 \omega+dd^c\psi=\sigma,
\] and the upper semicontinuous representative of (3) is \(\omega\)-psh
and nonpositive. It is smooth outside \(K\).

If \(x\in K\), use its level-\(j\) interval. On that interval
\(|x-w|\leq\delta_j\), while both denominator factors in \(G\) are at
least one. On the remainder of \(K\), \(G\leq0\). Thus \[
 \psi(x)\leq2^{-j}\log\delta_j=-2^j.
\] It follows that \(\psi=-\infty\) on \(K\). In particular, \(K\) is
polar and has area zero. The Lelong number of \(\psi\) at a point equals
the atom of its curvature measure there; all these Lelong numbers are
zero.

For \(0\leq\tau<1\), the ordinary curvature current of \(\tau\psi\) is
\[
 \omega+dd^c(\tau\psi)=(1-\tau)\omega+\tau\sigma.
\] In dimension one, its non-pluripolar part is its restriction off the
pole set. The whole measure \(\sigma\) is supported on that set.
Consequently, \[
 \int_{\mathbf P^1}\langle\omega+dd^c(\tau\psi)\rangle=1-\tau.
 \tag{4}
\] This computation is the reason for using a polar Cantor set; an
arbitrary Cantor set would not give (4).

\subsection{3. All exponential powers are
integrable}\label{all-exponential-powers-are-integrable}

For each \(b>0\), \[
 C_b:=\int_{\mathbf P^1}e^{-b\psi}\,\omega<\infty.
 \tag{5}
\] Here is a direct argument. At a point \(z_0\) in an affine chart,
choose \(r>0\) so small that \(q=\sigma(B(z_0,2r))<2/b\); this is
possible because \(\sigma\) has no atoms. For \(z\in B(z_0,r)\), the
kernel contribution from \(w\notin B(z_0,2r)\) is bounded below by a
finite constant. If \(q=0\) there is no near contribution. Otherwise,
Jensen's inequality for the normalized near measure gives \[
 \exp\left(-b\int_{B(z_0,2r)}G(z,w)\,d\sigma(w)\right)
 \leq\frac1q\int_{B(z_0,2r)}e^{-bqG(z,w)}\,d\sigma(w).
\] On the compact coordinate region, \(e^{-bqG(z,w)}\) is bounded by a
constant times \(|z-w|^{-bq}\). Its integral in \(z\) is uniformly
finite since \(bq<2\). Fubini gives local integrability of
\(e^{-b\psi}\). Away from the compact support, including infinity,
\(\psi\) is smooth. A finite cover proves (5). This proof provides a
separate finite constant for every fixed \(b\); no uniform bound as
\(b\to\infty\) is needed.

\subsection{4. A maximal test curve, with its endpoint
specified}\label{a-maximal-test-curve-with-its-endpoint-specified}

Let \[
 P_\omega[f]=\left(\lim_{C\to\infty}P_\omega(f+C,0)\right)^*,
\] where \(P_\omega(f+C,0)\) is the rooftop envelope. For \(0<\tau<a\),
put \(\chi_\tau=P_\omega[\tau\psi]\). Put \(\chi_\tau=0\) for
\(\tau\leq0\), set \[
 \chi_a=\lim_{\tau\uparrow a}\chi_\tau,
\] and set \(\chi_\tau=-\infty\) for \(\tau>a\).

We import the following positive-mass envelope theorem, with exactly its
usual hypotheses: if \(f\leq0\) is \(\omega\)-psh and its non-pluripolar
mass is positive, then \(P_\omega[f]\) is model, dominates \(f\), and
has the same non-pluripolar mass. This is the mass preservation and
idempotence established in Darvas--Di Nezza--Lu, Section 3.2 and Theorem
3.12. By (4), these hypotheses hold throughout \(0<\tau<a<1\). Therefore
\[
 \tau\psi\leq\chi_\tau\leq0,\quad
 P_\omega[\chi_\tau]=\chi_\tau,\quad
 \int\langle\omega+dd^c\chi_\tau\rangle=1-\tau.
 \tag{6}
\]

For completeness, the test-curve conditions are checked here.
Monotonicity follows from \(\psi\leq0\) and monotonicity of rooftop
envelopes. For concavity, take rooftop competitors below
\(\tau_1\psi+C_1\) and \(\tau_2\psi+C_2\), both below zero. Their convex
combination is an admissible competitor for \(\tau\psi+C\), where
\(\tau=\lambda\tau_1+(1-\lambda)\tau_2\) and
\(C=\lambda C_1+(1-\lambda)C_2\). Taking increasing envelope limits
gives \[
 \chi_\tau\geq\lambda\chi_{\tau_1}+(1-\lambda)\chi_{\tau_2}.
\] If one parameter is negative and the other positive, the same
argument uses \(0\leq\tau_1\psi\) for the negative parameter; when the
resulting parameter is nonpositive the inequality follows directly from
\(\chi_\tau=0\). Limits preserve concavity at \(a\); terms equal to
\(-\infty\) cause no extra condition.

For \(x\notin K\), the parameter function is finite and concave on
\((0,a)\), hence continuous there, and (6) gives continuity from the
right at zero. For any interior parameter \(\tau\), the decreasing left
limit is a psh function, since it is bounded below by \(\tau\psi\), and
agrees with \(\chi_\tau\) outside the area-zero set \(K\). Equality of
psh representatives from almost-everywhere equality gives equality
everywhere. Thus the parameter function is upper semicontinuous at every
interior parameter. The endpoint definition gives the same property at
\(a\); monotonicity handles right limits. At zero, upper semicontinuity
follows from \(\chi_\tau\leq0\) and the assigned value zero.

The endpoint is not identically \(-\infty\), since \(\chi_a\geq a\psi\).
Its model property requires no assumption about preservation of mass at
a zero-mass endpoint: for every \(\tau<a\), \[
 \chi_a\leq P_\omega[\chi_a]\leq P_\omega[\chi_\tau]=\chi_\tau.
\] Letting \(\tau\uparrow a\) gives \(P_\omega[\chi_a]=\chi_a\). This
proves that \(\chi\) is a bounded maximal test curve in the \textbf{psh
test-curve} terminology. It is not asserted to be an
\(\mathcal I\)-maximal test curve.

\subsection{5. The actual bounded weak geodesic ray and its
energy}\label{the-actual-bounded-weak-geodesic-ray-and-its-energy}

Use the inverse Legendre transform \[
 r_t=\sup_{\tau\in\mathbf R}(\chi_\tau+t\tau),\qquad
 v_t=r_t-at,\qquad t\geq0.
 \tag{7}
\] The Ross--Witt Nyström correspondence for bounded maximal test curves
supplies an actual bounded weak geodesic ray from zero. We use its form
in Darvas--Xia, Theorem 3.7(iii), with the energy formula from Theorem
3.7(iv) and Proposition 3.8. In the probability-volume convention it
states \[
 E_\omega(r_t)/t=\tau^+_\chi+
 \int_{-\infty}^{\tau^+_\chi}
 \left(\int\langle\omega+dd^c\chi_\tau\rangle-1\right)d\tau.
 \tag{8}
\] These are imported established results, not a claimed new proof of
the test-curve correspondence. Their boundedness hypothesis holds
because \(\chi_\tau=0\) for \(\tau\leq0\) and \(\tau^+_\chi=a\).
Equation (6) and (8) give \[
 E_\omega(r_t)=(a-a^2/2)t,\qquad E_\omega(v_t)=-a^2t/2.
 \tag{9}
\] Adding an affine function of time preserves the weak geodesic
equation. Since \(0\leq r_t\leq at\), one has \(-at\leq v_t\leq0\). The
term \(\tau=a\) in (7) also gives the stronger, time-independent
singular lower bound \[
 v_t\geq\chi_a\geq a\psi.
 \tag{10}
\] Every finite-time potential is bounded and therefore belongs to the
genuine finite-energy space. A fixed lower bound by the singular,
infinite-energy function \(a\psi\) does not make the ray constant.

\subsection{6. The entropy slope is infinite: a direct
estimate}\label{the-entropy-slope-is-infinite-a-direct-estimate}

For a probability measure \(\nu\) and reference probability measure
\(\mu\), relative entropy satisfies \[
 \operatorname{Ent}(\nu\mid\mu)\geq\int f\,d\nu-\log\int e^f\,d\mu
 \tag{11}
\] for every bounded Borel \(f\). If \(\nu\) is not absolutely
continuous the left side is infinite. Otherwise Jensen's inequality
applied to the density proves (11); truncation handles its zero set. In
particular (11) does not require the bounded potentials in (7) to be
continuous.

For bounded negative \(\omega\)-psh \(v\), the energy supporting-plane
inequality gives \[
 \int(-v)\,\operatorname{MA}_\omega(v)\geq-E_\omega(v).
 \tag{12}
\] On a curve this can also be seen directly: the difference between the
left side and the right side is \(-\tfrac12\int v\,dd^cv\geq0\), by
integration by parts. This identity is valid for bounded subharmonic
potentials by their usual Dirichlet-energy approximation.

Fix \(b>0\), and apply (11) with \(f=-bv_t\). By (5) and (10), \[
 \int e^{-bv_t}\omega\leq C_{ab}<\infty
\] uniformly in \(t\). Equations (9), (11) and (12) yield \[
 \operatorname{Ent}_\omega(v_t)\geq\frac{ba^2}{2}t-\log C_{ab}.
 \tag{13}
\] First let \(t\to\infty\) with \(b\) fixed, and then let
\(b\to\infty\). Thus the liminf of entropy divided by time is infinite,
and its extended-real limit is \(+\infty\). No interchange of
approximation and time limits, and no bound uniform in \(b\), is used.
The argument also allows the entropy at some finite times to be
infinite; finite energy of a bounded potential does not imply finite
entropy.

\subsection{7. The total ray has zero divisorial
singularities}\label{the-total-ray-has-zero-divisorial-singularities}

Let \(U(x,\zeta)=v_{-\log|\zeta|}(x)\) on \(\mathbf P^1\times D^*\). It
is \(p_1^*\omega\)-psh, is nonpositive, and extends as a psh potential
across the central divisor by upper semicontinuous extension. From (10),
for every \(b>0\), \[
 \int_{\mathbf P^1\times D}e^{-bU}\,\omega\wedge dA(\zeta)
 \leq\operatorname{Area}(D)\,C_{ab}<\infty.
 \tag{14}
\] Consider any proper modification \(\pi:Z\to\mathbf P^1\times D\) with
\(Z\) smooth, and a prime divisor \(F\) on \(Z\). Near a generic point
of \(F\), choose coordinates with \(F=\{z_1=0\}\), and write the
holomorphic Jacobian as \(z_1^k\) times a unit. If \(U\circ\pi\) had
generic Lelong number \(\lambda>0\) along \(F\), its divisorial Siu
decomposition would imply the local upper bound \[
 U\circ\pi\leq\lambda\log|z_1|+C.
\] The change-of-variables formula applied to (14) would then give
integrability of a function bounded below by a positive constant times
\(|z_1|^{-b\lambda+2k}\). This is impossible for \(b\lambda\geq2(k+1)\).
Since (14) holds for every \(b\), all such \(\lambda\) are zero.

This argument concerns divisors on the \textbf{total space}, including
divisors over the central fiber; it does not infer their vanishing
merely from zero Lelong numbers of finite-time slices. Under the stated
non-Archimedean convention every divisorial value of this ray is zero.
The usual divisorial determination of psh non-Archimedean metrics
identifies its datum with the constant metric zero. Alternatively, on
the projective curve one may check Darvas--Xia Proposition 3.13 directly
after shifting the test curve by \(a\): every nontrivial term has
divisorial Lelong number zero, and the supremum of its shifted
parameters is zero.

\subsection{8. Products and nonprojective
examples}\label{products-and-nonprojective-examples}

Define \(u_t=p_1^*v_t\) on \(X=\mathbf P^1\times Y\). The pullback of
the curve subgeodesic is a \(\Theta\)-subgeodesic. To check that it is
the weak geodesic between any two finite times, restrict an arbitrary
Perron competitor to \(\mathbf P^1\times\{y\}\) and its time annulus.
Its boundary values are those of the curve geodesic, so it lies below
\(v_t\) by the defining curve Perron property. Conversely, the pullback
of \(v_t\) is an admissible competitor. The two Perron envelopes
therefore agree. This establishes the actual product weak ray without
using a formal metric completion.

For a bounded curve potential \(v\) the Bedford--Taylor product gives \[
 \operatorname{MA}_\Theta(p_1^*v)
 =\operatorname{MA}_\omega(v)\otimes\frac{\eta^m}{\int_Y\eta^m}.
 \tag{15}
\] Indeed \(\Theta_{p_1^*v}=p_1^*\omega_v+p_2^*\eta\), and in its
\((m+1)\)st power the only nonzero term is
\((m+1)p_1^*\omega_v\wedge p_2^*\eta^m\). The same coefficient occurs in
\(V_X=(m+1)\int_Y\eta^m\) and cancels. Differentiating energy along the
segment \(s\mapsto p_1^*(sv)\) and using (15), with energy zero at
\(s=0\), gives \[
 E_\Theta(p_1^*v)=E_\omega(v).
 \tag{16}
\] This also follows from the defining energy sum. Equation (15) makes
the entropy exactly \(\operatorname{Ent}_\omega(v)\), including the case
of infinite entropy, because the second probability factor is the same
for the Monge--Ampère measure and the reference measure.

The total product ray is again bounded below by \(a\,p_1^*\psi\) and
above by zero. Fubini and (5) give local integrability of every
exponential power on \(X\times D\). The Jacobian argument from Section
7, now in dimension \(m+2\), still uses one transverse coordinate to a
divisor and proves zero generic divisorial Lelong numbers on every
modification. Hence the \textbf{actual} non-Archimedean datum on the
complex-space analytification is zero, rather than an arbitrarily
assigned datum. Mesquita--Piccione's Section 5.1 gives precisely this
total-ray construction and divisorial determination for compact Kähler
manifolds.

The normalized Monge--Ampère measure of the zero metric can be computed
on the trivial configuration \(X\times\mathbf P^1\): its sole central
component has multiplicity one, the restriction of the class is
\([\Theta]\), and its normalized top self-intersection is one. It thus
gives \(\delta_{v_{\rm triv}}\). With discrepancy over \(X\) normalized
to zero at the trivial valuation, (2) follows. This uses only the
defining model measure, not a general entropy-slope comparison theorem.

If \(Y\) is nonprojective then \(X\) is nonprojective, since a fiber
\(\{x\}\times Y\) would be a projective closed complex submanifold of a
projective \(X\). Such compact Kähler \(Y\) exist, for example
nonprojective complex tori of dimension two. Equations (9), (13), (15)
and (16) prove (1) in these genuinely nonpolarized examples as well.

\subsection{9. Consequences and the remaining original
question}\label{consequences-and-the-remaining-original-question}

The zero ray and \((u_t)\) have the same usual non-Archimedean datum,
but their entropy slopes are respectively \(0\) and \(+\infty\). Thus no
functional of that datum alone can equal the entropy slope for every
finite-energy weak ray. The failure already occurs for rays whose
finite-time potentials are bounded and whose non-Archimedean datum has
finite energy and entropy.

The term maximal in Section 4 describes \(P\)-maximal test curves, as
required for weak geodesics. The ray in (1) is not maximal relative to
its non-Archimedean datum: its maximal ray is the zero ray, while its
energy decreases with strictly negative slope. This distinction
preserves the hypotheses of the known projective maximal-ray theorems.
The example is consistent with Chi Li's finite-Mabuchi-slope maximality
theorem; its direct entropy proof above does not depend on that theorem.

Changing the smooth reference probability volume form only adds an
entropy term bounded uniformly in \(t\), since the logarithm of the
ratio of two smooth positive volumes is bounded on the compact manifold.
The infinite entropy slope therefore survives that change. Replacing
time by \(t' = c t\) for \(c>0\) multiplies the finite energy slope by
\(1/c\) and leaves the infinite entropy slope and zero non-Archimedean
datum unchanged.

The original P19.1 task asks for the residual entropy-slope statement
with its precise ray hypotheses. This candidate definitively rules out
the \textbf{unrestricted equality using the usual datum}. It does not
establish the missing upper entropy recovery for all maximal rays in
arbitrary transcendental Kähler classes. An independent reviewer must
determine whether a correctly formulated negative answer is sufficient
for the original atom; until then its solved count remains unchanged.

\subsection{10. Source applications and verification
scope}\label{source-applications-and-verification-scope}

Only three established inputs enter the proof: positive-mass
model-envelope preservation; the bounded-test-curve/weak-ray
correspondence and its normalized energy formula; and the standard
construction and divisorial determination of the non-Archimedean datum
of a Kähler ray. The Cantor measure, integrability estimates, endpoint
check, entropy divergence, total-space divisorial vanishing, product
Perron argument, and normalized product calculations are provided above.

\begin{itemize}
\tightlist
\item
  T. Darvas, E. Di Nezza, C. H. Lu, \emph{Monotonicity of non-pluripolar
  products and complex Monge--Ampère equations with prescribed
  singularity}, arXiv:1705.05796v3, Section 3.2 and Theorem 3.12. Root
  read the positive-mass mass-preservation discussion and the complete
  Theorem 3.12 proof. Its underlying monotonicity and convergence
  theorems are imported foundations.
\item
  T. Darvas, M. Xia, \emph{The closures of test configurations and
  algebraic singularity types}, Advances in Mathematics 397 (2022),
  108198, DOI 10.1016/j.aim.2022.108198, Theorem 3.7 and Proposition
  3.8. Root read the complete proof units, including Proposition 3.6,
  the full bounded-ray energy proof, and Lemma 3.9. Example 2.19
  motivates the construction, but (3)--(5) replace any unqualified
  assertion about a Cantor set. The local text is the published version
  of record, despite its cache filename containing the arXiv identifier
  2003.04818.
\item
  P. Mesquita--Piccione, \emph{A non-Archimedean theory of complex
  spaces and the cscK problem}, arXiv:2409.06221v2, Section 4.2.1 and
  Theorem 5.1.4 with Remark 5.1.5. Root read the model measure
  definition and the complete local datum construction and proof. Only
  the zero metric is used for the non-Archimedean measure calculation.
  The discrepancy over \(X\) is explicitly normalized above to avoid
  confusing it with the Gauss discrepancy over the total space.
\end{itemize}

No third-party source PDF or lengthy source extract is included in this
candidate. No native campaign input has been modified, no active
independent-review workspace has been inspected, and no feedback has
been sent to either active reviewer. At preparation, rounds 125 and 129
are the two running native Xean reviews using the requested gpt-6.1-sol
xhigh route. Original totals remain 7 solved and 35 open; independently
accepted scoped papers remain 49.

\end{document}
